% !TeX program = lualatex
% !TeX encoding = UTF-8
\documentclass[11pt,a4paper]{ltjsarticle}

\usepackage{amsmath,amssymb,mathtools}
\usepackage{amsthm}
\usepackage{lmodern}
\usepackage{booktabs}
\usepackage{longtable}
\usepackage{tabularx}
\usepackage{array}
\usepackage{enumitem}
\usepackage{geometry}
\usepackage{xcolor}
\usepackage{hyperref}

\geometry{margin=25mm}
\hypersetup{
  colorlinks=true,
  linkcolor=blue!55!black,
  citecolor=blue!55!black,
  urlcolor=blue!55!black,
  pdftitle={五句計算論},
  pdfauthor={千葉将臣}
}
\setlist{nosep}
\setlength{\emergencystretch}{3em}
\renewcommand{\arraystretch}{1.18}

\newtheorem{definition}{定義}[section]
\newtheorem{proposition}[definition]{命題}
\newtheorem{theorem}[definition]{定理}
\newtheorem{conjecture}[definition]{予想}
\newtheorem{corollary}[definition]{系}
\newtheorem{axiom}[definition]{公理}
\theoremstyle{remark}
\newtheorem*{remark}{注意}
\renewcommand{\proofname}{証明}
\renewcommand{\abstractname}{要旨}

\newcommand{\HoTT}{\mathrm{HoTT}}
\newcommand{\Maybe}{\mathsf{Maybe}}
\newcommand{\List}{\mathsf{List}}
\newcommand{\State}{\mathsf{State}}
\newcommand{\Cont}{\mathsf{Cont}}
\newcommand{\Id}{\mathrm{Id}}
\newcommand{\isContr}{\mathrm{isContr}}
\newcommand{\EM}{\mathrm{EM}}
\newcommand{\LawvereInfinity}{\Lambda_\infty}
\newcommand{\SelfRef}{\Sigma}
\newcommand{\Skobs}{\mathop{\mathrm{Sk}}_{\mathrm{obs}}}
\newcommand{\Skiso}{\mathop{\mathrm{Sk}}_{\mathrm{iso}}}

\title{五句計算論：記述可能性の段階としての\\
圏論・時間性・型理論・微分幾何学の統合}
\author{千葉将臣}
\date{2026-08-16}

\begin{document}

\maketitle

\begin{abstract}
本論文は、インド古代論理学のサンジャヤ五句否定
（Pañca-mūla-avaktavya）を「記述可能性の段階（stages of
describability）」として再解釈し、現代の圏論・ホモトピー型理論
（HoTT）・計算論・微分幾何学（diffeology）・様相論理を統合する包括的枠組み
「五句計算論（Penta-Computational Theory）」を提案する。五句分別の各句は、
型理論におけるh-level階層、圏論における随伴・モナド・コモナド、
微分幾何学における商空間・プロットの持ち上げ、様相論理における
必然演算子□と可能演算子◇の時間構造と精密に対応する。特に、HoTTのパス構造は
「時間の軌跡」として解釈され、truncationは「観測者の固定・証明の時間性の消去」
として理解される。本論文は、これらの対応を公理化し、未解決課題を明示する。

\medskip
\noindent\textbf{キーワード：}
五句分別、ホモトピー型理論、モナド、コモナド、様相論理、diffeology、
記述可能性、時間性、商型、truncation、Lawvere--Tierney topology、計算論
\end{abstract}

\section{序論：問題提起}

\subsection{五句分別の再解釈}

サンジャヤ学派の五句否定は、一つの命題 $P$ に対して以下の五つの
記述可能性を区別する。

\begin{enumerate}
  \item \textbf{$P$（構成的に真）}：証明 $p:P$ を持つ。
  \item \textbf{$\neg_i\neg_i P$（構成的になくない）}：間接的に到達可能である。
  \item \textbf{$\neg_c P$（非構成的にない）}：古典的否定（独立命題を含む）。
  \item \textbf{$\neg_i\neg_c P$（非構成的になくない）}：古典的に二重否定
  （証人を忘れた存在）。
  \item \textbf{矛盾でない（整合性）}：体系の無矛盾性。
\end{enumerate}

本論文の主張は、これらは「真理の度合い」ではなく、
\textbf{「記述可能性の段階（stage of describability）」}の区別である、
というものである。

\subsection{統合の動機}

現代数学・論理学には、以下の断絶がある。

\begin{itemize}
  \item \textbf{圏論}：構造の保存と段階間の移行（随伴）を記述するが、
  「何が記述されるか」の段階論を欠く。
  \item \textbf{ホモトピー型理論（HoTT）}：証明のパス構造を幾何学的に
  記述するが、時間性の哲学的解釈が分散している。
  \item \textbf{計算論}：モナドとしての計算効果を記述するが、その
  「段階的」構造が明示されていない。
  \item \textbf{微分幾何学（diffeology）}：商空間と滑らか構造を記述するが、
  論理的基礎との接続が未整備である。
  \item \textbf{様相論理}：$\Box/\Diamond$による「必然／可能」を記述するが、
  「世界」「到達可能性」などのメタファーが実装依存で曖昧である。
\end{itemize}

五句計算論は、これらを\textbf{「記述可能性の段階」という単一の軸}で統合する。

\section{五句計算論の公理化}

\subsection{基本定義}

\begin{definition}[記述可能性の段階]
命題 $P$ の記述可能性段階 $\delta(P)$ は、次の5値をとる。
\[
  \delta(P)\in\{S_1,S_2,S_3,S_4,S_5\}.
\]
各段階を表\ref{tab:stages}のように定義する。
\end{definition}

\begin{table}[htbp]
\centering
\caption{記述可能性の五段階}
\label{tab:stages}
\small
\begin{tabularx}{\textwidth}{@{}c l >{\raggedright\arraybackslash}X >{\raggedright\arraybackslash}X@{}}
\toprule
記号 & 名称 & 型理論的定義 & 直観 \\
\midrule
$S_1$ & 構成的に真 & $\exists p:P$（要素を持つ） & 今ここに証明がある \\
$S_2$ & 構成的になくない & $\neg\neg P$（二重否定） & 間接的に到達可能 \\
$S_3$ & 非構成的になくない & $\lVert P\rVert_{-1}$（命題切捨て） & 証人を忘れた存在 \\
$S_4$ & 非構成的にない & $\neg_cP$（古典的否定） & 体系内で証明不能 \\
$S_5$ & 矛盾でない & $\isContr(\lVert P\rVert_{-1}+\lVert\neg P\rVert_{-1})$ & 体系の整合性 \\
\bottomrule
\end{tabularx}
\end{table}

\subsection{h-level階層との対応}

HoTTのh-level階層は、五句分別を\textbf{無限に精密化}する
（表\ref{tab:hlevel}）。

\begin{table}[htbp]
\centering
\caption{h-level階層との対応}
\label{tab:hlevel}
\small
\begin{tabularx}{\textwidth}{@{}c >{\raggedright\arraybackslash}X >{\raggedright\arraybackslash}X@{}}
\toprule
h-level & 五句分別的読み & 具体例 \\
\midrule
$-2$ & $S_5$（矛盾でない／唯一の真理） & 可縮型 $\isContr(A)$ \\
$-1$ & $S_3$（非構成的になくない） & 命題 $\lVert A\rVert_{-1}$ \\
$0$ & $S_2$（構成的になくない／同一性証明が一意） & 集合 $\lVert A\rVert_0$ \\
$1$ & $S_1$の豊かさ（パスが複数） & 群胚 \\
$\geq 2$ & 記述可能性の完全展開 & 高次群胚、無限coherence \\
\bottomrule
\end{tabularx}
\end{table}

\begin{remark}
$S_2$（$\neg\neg P$）と$S_3$（$\lVert P\rVert_{-1}$）の区別は、
排中律（LEM）なしでのみ成立する。LEMのもとでは
$\lVert P\rVert_{-1}\simeq\neg\neg P$となるが、構成的文脈では両者は
\textbf{異なる段階}である。
\end{remark}

\section{Lawvere無限階層の五句計算論への統合}
\label{sec:lambda-infinity}

本節では$\LawvereInfinity$をLawvereの固定点定理の五句計算論的実装として、
圏論的言語のみで定式化する。対角補題、証明可能性述語、等号の二重性、
Priestの矛盾許容論理との差異、およびリーマン球面による解釈を含む
算術的・証明論的構成については、$\LawvereInfinity$単独論文
\cite{ChibaLambda2026}を参照されたい。

\subsection{五句段階における位置づけ}

\begin{definition}[Lawvere無限階層]
各五句段階$S_i$を表す圏を$\mathcal{C}_i$とする。
$\LawvereInfinity$は、$\mathcal{C}_4$（体系内不可記述）におけるLawvere固定点を、
五句段階循環
\[
 S_1\longrightarrow S_2\longrightarrow S_3\longrightarrow S_4
 \longrightarrow S_5\longrightarrow S_1
\]
の各周回に沿って集めた両立的な固定点族、またはその極限を表す記号である。
\end{definition}

\begin{table}[htbp]
\centering
\caption{$\LawvereInfinity$を構成する段階的構造}
\label{tab:lambda-stage}
\small
\begin{tabularx}{\textwidth}{@{}>{\raggedright\arraybackslash}X c >{\raggedright\arraybackslash}X@{}}
\toprule
構造 & 五句段階 & 役割 \\
\midrule
対角写像$\Delta_i:X_i\to X_i\times X_i$ & $S_1$--$S_5$ & 各段階での自己適用の舞台 \\
評価写像$\operatorname{eval}_i:X_i\times Y_i^{X_i}\to Y_i$ & $S_1$--$S_5$ & 指数対象の評価 \\
Lawvere固定点$b_i=g_i(b_i)$ & $S_4\to S_5$ & 各周回における穴の検出と閉包 \\
$\LawvereInfinity$ & $S_4$の無限階層化 & 5-periodic固定点列の極限的記述 \\
\bottomrule
\end{tabularx}
\end{table}

\subsection{段階的自己言及と5-periodic反復}

\begin{definition}[段階的Lawvereデータ]
各$\mathcal{C}_i$に対象$X_i,Y_i$と射
\[
 e_i:X_i\longrightarrow Y_i^{X_i}
\]
を置く。$e_i$が弱い点全射であるとは、任意の射$h:X_i\to Y_i$が、ある
$x:1\to X_i$に対する$e_i\circ x$の評価によって表現されることをいう。
自己写像$g_i:Y_i\to Y_i$へLawvereの対角構成を適用して得られる固定点を
$b_i:1\to Y_i$と書く。
\end{definition}

\begin{definition}[5-periodic反復列]
段階遷移関手$R_i:\mathcal{C}_i\to\mathcal{C}_{i+1}$（添字は5を法として読む）
がLawvereデータを保存すると仮定する。$X_0\in\mathcal{C}_1$に対して
\[
\begin{aligned}
X_1&=R_1(X_0)\in\mathcal{C}_2,&
X_2&=R_2(X_1)\in\mathcal{C}_3,\\
X_3&=R_3(X_2)\in\mathcal{C}_4,&
X_4&=R_4(X_3)\in\mathcal{C}_5,\\
X_5&=R_5(X_4)\in\mathcal{C}_1
\end{aligned}
\]
を一周とする。各周回の$\mathcal{C}_4$成分に得られる固定点$b_4^{(k)}$が
遷移関手と両立するとき、これらは固定点塔をなす。
\end{definition}

\begin{definition}[$\LawvereInfinity$の極限表示]
$\mathcal{C}_4$における固定点塔$(b_4^{(k)})_{k\geq0}$の極限が存在するとき、
\[
 \LawvereInfinity:=\varprojlim_k b_4^{(k)}
\]
と定義する。各成分は対応する自己写像$g_4^{(k)}$に対して
$g_4^{(k)}(b_4^{(k)})=b_4^{(k)}$を満たす。
\end{definition}

\begin{remark}[等号の意味]
上式の等号は、$A=A$という対象の反射的自己同一性と区別される。
ここでは$\lim_{n\to\infty}a_n=L$や$x=f(x)$と同様に、固定点塔の極限を
静的に表示する。したがって、この等号は$S_3$から直接観測できない$S_4$の構造へ
名前を与える圏論的な極限表示である。等号の証明論的二重性は
\cite{ChibaLambda2026}で扱う。
\end{remark}

\subsection{存在条件と五句公理の拡張}

\begin{axiom}[Lawvere点全射条件]
$\mathcal{C}_4$における射$e_4:X_4\to Y_4^{X_4}$が、Lawvereの固定点定理に
必要な弱い点全射条件を満たすと仮定する。
\end{axiom}

\begin{theorem}[$\LawvereInfinity$の内在的存在]
五句公理C1--C5、$S_3\to S_4$の非可逆な段階移行、Lawvere点全射条件、
および固定点塔の極限存在を
満たす五句計算体系$\mathcal{K}$では、
\[
 \LawvereInfinity=\varprojlim_k b_4^{(k)}\in\mathcal{C}_4
\]
が成立する。
\end{theorem}

\begin{proof}[証明の概略]
$S_3\to S_4$の非可逆性は、観測可能な商から失われる構造を生じさせる。
ただし多対一性だけでは固定点の存在は従わないため、上の点全射条件を明示的に用いる。
Lawvereの固定点定理を$\mathcal{C}_4$の対角写像と評価写像へ適用すると、
各周回の自己写像$g_4^{(k)}$に固定点$b_4^{(k)}$が得られる。遷移関手との両立性に
より固定点塔が形成され、仮定した極限の普遍性から$\LawvereInfinity$が得られる。
\end{proof}

\begin{axiom}[C4-$\Lambda$：$S_4$の自己言及公理]
$\mathcal{C}_4$は弱い点全射なLawvereデータと、その固定点塔を内在的に含む。
\[
 \LawvereInfinity=\varprojlim_k b_4^{(k)}\in\mathcal{C}_4.
\]
\end{axiom}

\begin{axiom}[C5-$\Lambda$：統制固定点との接続]
$S_5$の統制固定点$T\cong F(T)$は、$S_4$の穴として現れる
$\LawvereInfinity$を閉包へ統合し、次の周回$S'_1$への入口を生成する。
\end{axiom}

\subsection{既存構造との統合}

\begin{table}[htbp]
\centering
\caption{$\LawvereInfinity$と五句計算論の既存構造}
\label{tab:lambda-integration}
\small
\begin{tabularx}{\textwidth}{@{}>{\raggedright\arraybackslash}p{0.29\textwidth} >{\raggedright\arraybackslash}X@{}}
\toprule
既存構造 & $\LawvereInfinity$との関係 \\
\midrule
$S_4$ & 体系内不可記述を具体化する自己言及固定点 \\
フロベニウス条件$\sim$ & 非可逆性が固定点構成の背景を与えるが、存在には点全射条件を要する \\
二諦公理 & $\LawvereInfinity$は$\Skobs$には現れず、$\Skiso$の外側を標示する \\
バナッハ＝タルスキー接続 & 多対一の観測圧縮が残す不可視構造の$S_4$レベルでの表現 \\
Frege--Lawvere接続 & SinnがBedeutungへ完全には潰れないことの固定点表示 \\
計算の本質$S_1\to S_3$ & 計算の外側$S_4$における自己言及であり、継続領域の圏論的候補 \\
\bottomrule
\end{tabularx}
\end{table}

\section{圏論的基礎：随伴・モナド・コモナド}

\subsection{随伴＝段階間の移行}

Lawvereの\textit{Adjointness in Foundations}\cite{Lawvere1967}によれば、
論理的基本操作（$\exists,\forall,\Rightarrow,\times,+$）は随伴として記述される。
Rodin\cite{Rodin2013}の解釈によれば、これはヘーゲル的弁証法（正・反・合）の
圏論的実現である。

五句計算論では、\textbf{随伴は「記述可能性の段階間の移行」}として再解釈される。
\[
  F\dashv G:\mathcal{C}\longrightarrow\mathcal{D}.
\]
ここで、$F$は段階の「持ち上げ」（自由化）、$G$は段階の「忘却」（制限）を表す。

\subsection{モナド＝段階の操作}

Moggi\cite{Moggi1991}およびPlotkin--Power\cite{PlotkinPower2002}の
「計算＝モナド」対応を五句分別軸で再解釈すると、表\ref{tab:monads}を得る。

\begin{table}[htbp]
\centering
\caption{モナドと五句分別の対応}
\label{tab:monads}
\small
\begin{tabularx}{\textwidth}{@{}c l >{\raggedright\arraybackslash}X@{}}
\toprule
モナド $T$ & 計算概念 & 五句分別対応 \\
\midrule
$\Maybe$ & 例外／部分性 & $S_1\to S_2$（失敗の可能性） \\
$\List$ & 非決定性 & $S_1$の証人の多重化 \\
$\State$ & 文脈依存 & $S_1$の時間的持続性 \\
$\Cont$ & 二重否定 & $S_2$の直接実現（$\neg\neg A$） \\
$\lVert\mathord{\cdot}\rVert_{-1}$ & 証人の忘却 & $S_1\to S_3$（段階上昇） \\
\bottomrule
\end{tabularx}
\end{table}

\begin{definition}[五句モナド]
モナド$(T,\eta,\mu)$が\textbf{五句モナド}であるとは、以下を満たすことをいう。
\begin{enumerate}
  \item $\eta:\Id\Rightarrow T$は$S_1\to S_{\geq 2}$の段階上昇である。
  \item $\mu:T^2\Rightarrow T$は「二重操作の単純化」、すなわち段階の準冪等性である。
  \item Beck条件（余等化子保存）により、段階の「忠実性」が保証される。
\end{enumerate}
\end{definition}

\subsection{コモナド＝必然性の固定}

様相論理的$\Box$はコモナド$W$として記述される（表\ref{tab:comonad}）。

\begin{table}[htbp]
\centering
\caption{コモナドと必然性}
\label{tab:comonad}
\small
\begin{tabularx}{\textwidth}{@{}>{\raggedright\arraybackslash}X >{\raggedright\arraybackslash}X >{\raggedright\arraybackslash}X@{}}
\toprule
公理 & コモナド則 & 五句計算論的解釈 \\
\midrule
$\Box P\to P$（T） & 余単位 $\varepsilon:W\Rightarrow\Id$ & 固定の解除＝現在へ回帰 \\
$\Box P\to\Box\Box P$（4） & 余乗法 $\delta:W\Rightarrow W^2$ & 固定の不変性 \\
$\Box\Box P\to\Box P$ & 余乗法のcounit性 & 二重固定の単純化 \\
\bottomrule
\end{tabularx}
\end{table}

\section{時間性の定式化}

\subsection{証明＝パス＝時間の軌跡}

HoTTにおいて、同一性証明$p:x=_Ay$はパス$p:\mathbb{I}\to A$として解釈される。
\begin{itemize}
  \item $\operatorname{refl}_a:a=_Aa$： 「今」の定数パス。
  \item $p\mathbin{\cdot}q:x=_Az$：未来の合成。
  \item $p^{-1}:y=_Ax$：過去の逆パス。
\end{itemize}

\begin{definition}[時間的一貫性]
Path induction（J rule）は、次の時間的一貫性の型理論的定式化である。
\[
J:\prod_{C:\prod_{y:A}(x=y)\to\mathcal{U}}
\left(C(x,\operatorname{refl}_x)\to
\prod_{y:A}\prod_{p:x=y}C(y,p)\right).
\]
これは「今の真理を未来に運ぶ」操作として解釈される。
\end{definition}

\subsection{観測者の固定＝truncation＝時間性の消去}

古典数学的証明の「非時間性（atemporality）」は、propositional truncation
$\lVert P\rVert_{-1}$による証明パスの同一視として記述される。
\[
 |p|:\lVert P\rVert_{-1},\quad |q|:\lVert P\rVert_{-1}
 \quad\Longrightarrow\quad
 |p|=_{\lVert P\rVert_{-1}}|q|.
\]

\begin{proposition}
古典数学的証明は、五句計算論的に$S_1\to S_3$のtruncation操作によって
得られる「時間性＝0」の対象である。
\end{proposition}

\begin{proof}[証明の概略]
構成的証明$p:P$（$S_1$）から命題切捨て$\lVert P\rVert_{-1}$（$S_3$）への
移行は、証明パス$p,q$の差異（時間性）を同一視する。これはproof irrelevanceの
型理論的実現であり、「観測者（証明者・証明パス）の差異を固定・忘却する」操作に
対応する。
\end{proof}

\section{HoTTとの接続}

\subsection{パス構造と無限coherence}

モナド則の厳密等式
$\mu\circ T\mu=\mu\circ\mu_T$は、HoTT的にはホモトピー
$\mu\circ T\mu\sim\mu\circ\mu_T$として緩和される。その上に2-cell、3-cell、
さらに高次のセルからなる\textbf{無限coherence tower}が構築される
\cite{RiehlVerity2015}。

\begin{proposition}[coherence tower＝時間の解像度の精密化]
モナドの無限coherence towerは、五句分別の$S_1\to S_5$の遷移を
\textbf{無限に分解}する操作である。各h-level $n$は、時間の解像度の$n$次の
微分に対応する。
\end{proposition}

\subsection{HITによる商型とdiffeology・EM圏の統合}

Higher Inductive Type（HIT）による商型$\operatorname{Quotient}(A,R)$は、
以下のconstructorを持つ。
\begin{itemize}
  \item point constructor：$|a|:\operatorname{Quotient}(A,R)$。
  \item path constructor：$\operatorname{eq}(a,b,r):|a|=|b|$（$r:R(a,b)$）。
\end{itemize}

これはdiffeologyの商空間$X\to X/G$およびEM圏の商代数
$T^2A\rightrightarrows TA\to A$（余等化子）と形式的に対応する。

\begin{theorem}[van Doorn]
HITsは商型に帰着可能である（HITs reduce to quotients）\cite{vanDoorn2016}。
\end{theorem}

\begin{corollary}
diffeologyの商空間、HoTTのHIT、EM圏の商代数は、五句計算論的に
\textbf{同一の段階操作}（$S_1\to S_3$の商同一視）の異なる実装である。
\end{corollary}

\subsection{Univalence＝米田的真理保存性}

Univalence Axiom
\[
 (A=_{\mathcal{U}}B)\simeq(A\simeq B)
\]
は、\textbf{「等価な構造は同一の記述可能性段階に属する」}という
五句計算論的原理の型理論的実現である\cite{HoTTBook2013}。

\section{diffeologyとの接続}

\subsection{商空間とプロットの持ち上げ}

diffeologyの基本構造は、集合$X$上の\textbf{プロット（plot）}の族
$dX\subset\{\mathbb{R}^n\to X\}$である\cite{IglesiasZemmour}。

\begin{table}[htbp]
\centering
\caption{diffeology・五句計算論・圏論の対応}
\label{tab:diffeology}
\small
\begin{tabularx}{\textwidth}{@{}>{\raggedright\arraybackslash}X >{\raggedright\arraybackslash}X >{\raggedright\arraybackslash}X@{}}
\toprule
diffeology & 五句計算論 & 圏論 \\
\midrule
自由プロット & $S_1$（構成的に真） & クライスリ圏$\mathcal{C}_T$ \\
任意プロット（商の後） & $S_3$（非構成的になくない） & EM圏$\mathcal{C}^T$ \\
D-topology & $S_3$の位相的実現 & 比較関手$K:\mathcal{C}_T\to\mathcal{C}^T$のfaithfulness \\
不合理トーラス$\mathbb{T}_\theta$ & $S_4$（非構成的にない） & 記述可能性の境界 \\
\bottomrule
\end{tabularx}
\end{table}

\subsection{Beck条件とsubductionの同値性（予想）}

\begin{conjecture}
モナド$T$がBeck条件（余等化子保存・忠実性）を満たすことと、
diffeology的subduction（商写像の持ち上げ性質）が同値である。
\end{conjecture}

Beck条件は「段階の移行が情報を失わない」ことを保証し、subductionは
「商空間の滑らか構造が原空間から持ち上げられる」ことを保証する。
両者は五句計算論的に\textbf{「$S_1\to S_3$の移行の可逆性条件」}として
同一視される。

\section{様相論理との統合}

\subsection{\texorpdfstring{$\Box/\Diamond$}{必然・可能演算子}の五句計算論的分解}

様相論理の$\Box$（必然）と$\Diamond$（可能）は、
\textbf{「記述可能性の段階を操作する普遍的演算子」}として再解釈される。

\begin{table}[htbp]
\centering
\caption{様相演算子の五句計算論的分解}
\label{tab:modal-op}
\small
\begin{tabularx}{\textwidth}{@{}c >{\raggedright\arraybackslash}X >{\raggedright\arraybackslash}X@{}}
\toprule
演算子 & 圏論的対応 & 五句計算論的解釈 \\
\midrule
$\Box P$ & コモナド$W$ & 記述可能性の「固定・完備化」 \\
$\Diamond P$ & モナド$M$ & 記述可能性の「拡張・持ち上げ」 \\
$\Box\Box P\to\Box P$ & 余乗法のcounit & 二重固定の単純化＝段階の準冪等性 \\
$\Diamond\Diamond P\to\Diamond P$ & 乗法$\mu$ & 二重拡張の平坦化 \\
\bottomrule
\end{tabularx}
\end{table}

\subsection{S4/S5の公理の五句計算論的読み}

\begin{table}[htbp]
\centering
\caption{S4/S5公理の解釈}
\label{tab:s4s5}
\small
\begin{tabularx}{\textwidth}{@{}>{\raggedright\arraybackslash}X >{\raggedright\arraybackslash}X@{}}
\toprule
公理 & 五句計算論的解釈 \\
\midrule
$\Box P\to P$（T） & 普遍的真理は現在の真理＝最上位段階からの回帰 \\
$\Box P\to\Box\Box P$（4） & 固定の不変性＝段階の自己保存 \\
$P\to\Box\Diamond P$（B） & 現在の真理は必然的に可能＝現在の記述はすべての未来から参照可能 \\
$\Diamond\Box P\to\Box P$（5） & 可能性の必然性は必然性＝段階の最上位収束＝EM圏の完備性 \\
\bottomrule
\end{tabularx}
\end{table}

\subsection{「世界」メタファーの解消：時間差分の導入}

Kripke意味論の「世界$w$」は、表\ref{tab:worlds}の三つの異なる概念を
混同している。

\begin{table}[htbp]
\centering
\caption{「世界」の三分類}
\label{tab:worlds}
\small
\begin{tabularx}{\textwidth}{@{}>{\raggedright\arraybackslash}X c >{\raggedright\arraybackslash}X@{}}
\toprule
「世界」の型 & 時間差分$\Delta t$ & 五句分別 \\
\midrule
同時の反事実的可能性 & $\Delta t=0$ & $S_2$（$\neg\neg P$） \\
分岐的未来 & $\Delta t>0$ & $S_1$のパス構造 \\
形而上学的可能性 & $\Delta t=\mathrm{undefined}$ & $S_4$（体系外部） \\
\bottomrule
\end{tabularx}
\end{table}

\begin{proposition}
様相論理の意味論が曖昧であるのは、「観測者を固定する」操作が
$\Delta t=0$の仮定（同時の別可能性）のみを扱い、$\Delta t>0$（分岐的未来）と
$\Delta t=\mathrm{undefined}$（形而上学）を区別しないためである。
\end{proposition}

五句計算論は、この区別を\textbf{h-levelの時間的拡張}として精密化する。

\section{計算論的解釈}

\subsection{Moggi--Plotkin--Powerの対応の再構成}

Moggi\cite{Moggi1991}およびPlotkin--Power\cite{PlotkinPower2002}によれば、
計算概念はモナドを決定する。五句計算論では、これは
\textbf{「計算効果＝記述可能性段階の操作」}として再解釈される。

\begin{table}[htbp]
\centering
\caption{計算概念と段階操作}
\label{tab:computation}
\small
\begin{tabularx}{\textwidth}{@{}l c >{\raggedright\arraybackslash}X@{}}
\toprule
計算概念 & モナド & 五句分別的段階操作 \\
\midrule
例外 & $\Maybe$ & $S_1\to S_2$（失敗の可能性の導入） \\
非決定性 & $\List$ & $S_1$の証人の多重化 \\
状態 & $\State$ & $S_1$の時間的持続性（$\Delta t>0$のパス） \\
継続 & $\Cont$ & $S_2$（$\neg\neg A$）の直接実現 \\
証人忘却 & $\lVert\mathord{\cdot}\rVert_{-1}$ & $S_1\to S_3$（時間性の消去） \\
\bottomrule
\end{tabularx}
\end{table}

\subsection{Lawvereの「随伴の基礎」の五句計算論的再構成}

Lawvere\cite{Lawvere1967}の命題「すべての論理的基本操作は随伴である」は、
五句計算論的に次のように再解釈される。
\begin{quote}
\textbf{段階間の移行そのものが随伴として実現される。}
\end{quote}
正（$S_1$）・反（$S_4$）・合（$S_3$）の弁証法は、随伴の合成
（モナド／コモナド）として記述される。

\section{統合図式}

\subsection{五句計算論の対応表（総合）}

\begin{table}[htbp]
\centering
\caption{五句計算論の総合対応表}
\label{tab:integrated}
\scriptsize
\setlength{\tabcolsep}{3pt}
\begin{tabularx}{\textwidth}{@{}c c >{\raggedright\arraybackslash}X >{\raggedright\arraybackslash}X >{\raggedright\arraybackslash}X >{\raggedright\arraybackslash}X >{\raggedright\arraybackslash}X@{}}
\toprule
五句 & h-level & HoTT & 圏論 & diffeology & 様相論理 & 計算論 \\
\midrule
$S_1$ & $\geq1$ & $p:P$（パス） & クライスリ圏 & 自由プロット & $\Diamond P$（$\Delta t>0$） & 純粋値 \\
$S_2$ & $0$ & $\neg\neg P$ & 随伴の合成 & プロットの持ち上げ & $\Diamond P$（$\Delta t=0$） & 間接到達 \\
$S_3$ & $-1$ & $\lVert P\rVert_{-1}$（HIT） & EM圏・商代数 & 任意プロット・D-topology & $\Box P$（$\Delta t=0$） & 証人忘却 \\
$S_4$ & --- & 独立命題 & 比較関手の非faithfulness & 不合理トーラス & $\Box\neg P$ & 計算不能 \\
$S_5$ & $-2$ & $\isContr$ & 圏の整合性 & 空のdiffeology & 体系の完全性 & 停止問題 \\
\bottomrule
\end{tabularx}
\end{table}

\subsection{三角対応（HoTT・diffeology・EM圏）}

\begin{figure}[htbp]
\centering
\[
\begin{array}{c}
\text{HoTT（HIT）}\quad \operatorname{Quotient}(A,R)\\[4pt]
\swarrow\quad\text{商同一視 }(S_1\to S_3)\quad\searrow\\[4pt]
\text{diffeology（プロット）}\quad\longleftrightarrow\quad
\text{EM圏（$T$-代数）}\\[4pt]
\text{プロットの持ち上げ}\quad=\quad
\text{比較関手 $K$}\quad=\quad\text{induction principle}
\end{array}
\]
\caption{HoTT・diffeology・EM圏の三角対応}
\label{fig:triangle}
\end{figure}

\section{未解決課題と今後の展望}

\subsection{形式的未解決項目}

\begin{table}[htbp]
\centering
\caption{形式的未解決項目（1--6）}
\label{tab:open-problems}
\small
\begin{tabularx}{\textwidth}{@{}c >{\raggedright\arraybackslash}X l@{}}
\toprule
番号 & 課題 & 優先度 \\
\midrule
1 & HoTT・LT位相・五句分別の形式対応
（$\lVert A\rVert$ vs. $\neg\neg A$ vs. $\Diamond A$） & 最高 \\
2 & Beck条件とdiffeology的subductionの同値性の証明 & 最高 \\
3 & クライスリ圏の4モナド（Maybe, List, State, Cont）のdiffeology/HoTT翻訳 & 最高 \\
4 & HoTTでのモナド無限coherenceの具体的型項構築 & 高 \\
5 & CLTL（計算樹状時相論理）との形式接続 & 高 \\
6 & Berger--Melliès--WeberのNerve theoremと五句分別の定理的対応 & 中 \\
\bottomrule
\end{tabularx}
\end{table}

\begin{table}[htbp]
\centering
\caption{形式的未解決項目（7--11）}
\small
\begin{tabularx}{\textwidth}{@{}c >{\raggedright\arraybackslash}X l@{}}
\toprule
番号 & 課題 & 優先度 \\
\midrule
7 & 実効的トポス（Hyland）と構成的側面の接続 & 中 \\
8 & D-topologyとLT位相の同型／忠実埋め込みの証明 & 中 \\
9 & 20状態$j_m(K_n)$の代数的・動的記述 & 低 \\
10 & 様相論理S5への完全翻訳（五句計算論的語彙での公理化） & 低 \\
11 & $S_4$におけるLawvere点全射条件と$\LawvereInfinity$存在条件のモデル構成 & 最高 \\
\bottomrule
\end{tabularx}
\end{table}

\subsection{哲学的未解決項目}

\begin{enumerate}
  \item \textbf{「観測者」の形式的定義}：truncationは「観測者の忘却」として
  解釈されるが、「観測者」そのものを型理論的に定式化する必要がある。
  \item \textbf{時間矢の非対称性}：コモナド（$\Box$）とモナド（$\Diamond$）の
  非対称性が、熱力学的時間矢とどのように関係するか。
  \item \textbf{第五句の幾何学的実現}：$S_5$（整合性）はh-level $-2$
  （可縮型）に対応するが、これが「体系の完全性」とどのように幾何学的に対応するか。
\end{enumerate}

\section{結論}

五句計算論は、サンジャヤ五句否定を「記述可能性の段階」として再解釈し、
圏論・HoTT・diffeology・様相論理・計算論を統合する枠組みである。

本論文の主たる貢献は以下のとおりである。
\begin{enumerate}
  \item \textbf{五句分別の型理論的精密化}：h-level階層との対応により、
  5段階から無限連続体への精密化を実現した。
  \item \textbf{モナド／コモナドの五句計算論的再解釈}：$\Box/\Diamond$、
  計算効果、truncationを「段階の操作」として統一した。
  \item \textbf{時間性の定式化}：証明＝パス＝時間の軌跡、
  truncation＝時間性の消去＝観測者の固定として定式化した。
  \item \textbf{様相論理の意味論の解消}：「世界」メタファーを
  「記述可能性の段階の文脈」に置き換え、$\Delta t$（時間差分）による
  三重分解を提示した。
  \item \textbf{HoTT・diffeology・EM圏の三角統合}：商型（HIT）、商空間
  （diffeology）、商代数（EM圏）の形式的対応を提示した。
  \item \textbf{Lawvere無限階層の統合}：$\LawvereInfinity$を$S_4$の
  Lawvere固定点塔の極限として定義し、その存在に必要な点全射条件を明示した。
  真理保存体系における算術的固定点は、この圏論的一般構造の具体例として分離した。
\end{enumerate}

五句計算論は、単なる哲学的類比ではなく、
\textbf{具体的な数学的対応（型・圏・幾何学）を伴う研究プログラム}として
提示される。

\begin{thebibliography}{99}

\bibitem{Lawvere1967}
F. W. Lawvere,
``Adjointness in Foundations,''
\textit{Reprints in Theory and Applications of Categories}, no.~16
(original work published 1967).

\bibitem{Moggi1991}
E. Moggi,
``Notions of Computation and Monads,''
\textit{Information and Computation}, vol.~93, no.~1, pp.~55--92, 1991.

\bibitem{PlotkinPower2002}
G. Plotkin and J. Power,
``Notions of Computation Determine Monads,''
in \textit{Foundations of Software Science and Computation Structures}, 2002.

\bibitem{RiehlVerity2015}
E. Riehl and D. Verity,
``Homotopy Coherent Adjunctions and the Formal Theory of Monads,''
\textit{Advances in Mathematics}, 2015.

\bibitem{vanDoorn2016}
F. van Doorn,
\textit{Constructing the Propositional Truncation using Non-recursive HITs},
2016.

\bibitem{HoTTBook2013}
The Univalent Foundations Program,
\textit{Homotopy Type Theory: Univalent Foundations of Mathematics},
Institute for Advanced Study, 2013.

\bibitem{IglesiasZemmour}
P. Iglesias-Zemmour,
\textit{Diffeology},
American Mathematical Society, 2013.

\bibitem{PfenningDavies2001}
F. Pfenning and R. Davies,
``A Judgmental Reconstruction of Modal Logic,''
\textit{Mathematical Structures in Computer Science}, vol.~11, no.~4, 2001.

\bibitem{dePaivaRitter2021}
V. de Paiva and E. Ritter,
``Doctrines, Modalities and Comonads,''
\textit{Mathematical Structures in Computer Science}, 2021.

\bibitem{Rodin2013}
A. Rodin,
``Categorical Logic and Hegelian Dialectics,''
\textit{Philomatica}, 2013.

\bibitem{Kavvos2024}
G. A. Kavvos,
``Two-dimensional Kripke Semantics I: Presheaves,''
in \textit{FSCD}, 2024.

\bibitem{Hyland}
J. M. E. Hyland,
``The Effective Topos,''
in \textit{The L. E. J. Brouwer Centenary Symposium}, 1982.

\bibitem{SEP}
Stanford Encyclopedia of Philosophy,
entries ``Branching Time,'' ``Temporal Logic,'' ``Modal Logic,''
``Possible Worlds,'' and ``Two-Dimensional Semantics.''

\bibitem{nLab}
nLab,
entries ``modal type theory,'' ``Lawvere--Tierney topology,''
``double negation,'' ``BHK interpretation,'' ``identity type,'' and ``interval type.''

\bibitem{ChibaLambda2026}
千葉将臣，
『真理保存性を持つ体系が必然的に$\LawvereInfinity$を含む：
対角補題によるLawvere無限階層の内在的定式化』，2026．

\end{thebibliography}

\end{document}